\begin{translatedabstract}
Understanding health-related technical guidance can be challenging due to specialized vocabulary, long sentences, and limited adaptation to the communicational profile of end users in mHealth applications. This work aims to investigate, implement, and evaluate a solution based on Large Language Models for the simplification and communicational adaptation of health-related technical texts, while preserving the essential content of the original material. To this end, SimplyHealth was developed and contextualized within Takere, a no-code platform for mHealth applications based on care plans, treated as a reference environment without full integration. The solution receives free-text input or PDF files containing health guidance, considers simulated patient profiles, and generates simplified versions, a glossary, frequently asked questions, and a contextual chat. The adaptation is strictly communicational, not clinical, using aspects such as age, education level, and field of study only to adjust vocabulary, tone, level of explanation, and presentation style. The implementation includes prompts with content-preservation rules, two-layer guardrails — deterministic checking and LLM-based evaluation — to mitigate meaning changes, and the Flesch Readability Index adapted for Brazilian Portuguese (Flesch-PT) as an automatic readability indicator. The functional and linguistic evaluation used two health materials, three simulated profiles, and 18 main executions. The guardrails acted in 8 out of 18 executions, blocking simplifications with identified problematic alterations. The readability results showed an increase in Flesch-PT in 15 out of 18 cases, with an average score of 46.8 for the original texts and 58.2 for the simplified versions, resulting in an average gain of 11.4 points. These findings indicate the functional and linguistic feasibility of the approach within the evaluated scope, without representing clinical validation or evidence of improved understanding by real patients.
\end{translatedabstract}
